\documentclass[10pt]{article}

\usepackage[utf8]{inputenc}

\usepackage[T1]{fontenc}

\usepackage{amsmath}

\usepackage{amsfonts}

\usepackage{amssymb}

\usepackage[version=4]{mhchem}

\usepackage{stmaryrd}

\usepackage{bbold}

\usepackage{graphicx}

\usepackage[export]{adjustbox}

\graphicspath{ {./images/} }



\begin{document}

Р. есть циклоида; если окружность катится по окружности - циклоидальная кривая; если по прямой катится гипербола, эллипс или парабола - III турма криеая. Траектория точки әллипса, катящегося по др. әллипсу, наз. э п и э л л и п с о м. Каждая плоская кривая многими способами может быть рассмотрена как Р.; напр., всякая кривая может быть образована качением прямой по әволюте.



Лит.: [1] С а в е л о в А. А., Плоские кривые, М., 1960.



РУНГЕ ОБЛАСТЬ, о б л а с т ь Р ун г е П е р в ог о р о д а, - область $G$ в пространстве $\mathbb{C}^{n}$ комплексных переменных $\left(z_{1}, \ldots, z_{n}\right)$, обладающая тем свойством, что для любой голоморфной в $G$ функции $f\left(z_{1}, \ldots, z_{n}\right)$ существует последовательность многочленов



\[

\left\{P_{k}\left(z_{1}, \ldots, z_{n}\right)\right\}_{k=1}^{\infty},

\]



сходящаяся в $G$ к $f\left(z_{1}, \ldots, z_{n}\right)$ равномерно на каждом замкнутом ограниченном множестве $E \subset G$. Определение Р. о. в т о р о г о р о д а получается отсюда заменой последовательности (*) последовательностью рациональных функций $\left\{R_{k}\left(z_{1}, \ldots, z_{n}\right)\right\}_{k=1}^{\infty}$. При $n=1$ всякая односвязная область является Р. о. первого рода, всякая область - Р. о. второго рода (см. Рунге теорема). При $n \geqslant 2$ не всякая односвязная область есть Р. о. и не всякая Р. о. односвязна.



Лит.: [1] Ф у к с Б. А , Специальные главы теории аналитических функций многих комплексных переменных, М., 1963; [2] В л а д и м и р о в В. С, Методы теории функций многих комплексных переменных, М., 1964.



РУНГЕ ПРАВИЛО - один из методов оценки погрешности формул численного интегрирования. Пусть $R=h^{k} M$ - остаточный член формулы численного интегрирования, где $h$ - длина отрезка интегрирования или какой-то его части, $k$ - фиксированное число и $M-$ произведение постоянной на производную подинтегральной функции порядка $k-1$ в какой-то точке промежутка интегрирования. Если $J$ - точное значение интеграла, а $I$ - его приближенное значение, то $J=$ $=I+h^{k} M$.



Согласно Р. п. вычисляется тот же самый интеграл по той же формуле численного интегрирования, но вместо $h$ берется величина $h / 2$. При этом, чтобы получить значение интеграла по всему отрезку, формула интегрирования применяется дважды. Если производная, входящая в $M$, меняется не сильно на рассматриваемом промежутке, то



\[

R=h^{k} M=\frac{I_{1}-I}{1-\frac{1}{2^{k-1}}},

\]



где $I_{1}$ - значение интеграла, вычисленное по $h / 2$. Р. п. используется и при численном решении дифференциальных уравнений.



Правило предложено К. Рунге (С. Runge, нач. 20 в.).



Jum.: [1] Б е р е з и н И. С., Ж и д к о в Н. П., Методы вычислений, 3 изд., т. 1, М., 1966; 2 изд., т. 2 , М., 1962; [2] Современные численные методы решения обыкновенных дифференциальных уравнений, пер. с англ., М., 1979.



PУНГE TEOPEMA - Tеорема А. Б. Иванов. номиальных приближений голоморфных функций, впервые доказанная К. Рунге (С. Runge, 1885).



Пусть $D$ - односвязная область на плоскости комплексного переменного z. Тогда всякая функция $f$, голоморфная в $D$, приближается равномерно внутри $D$ многочленами от $z$. Точнее, для любого компакта $K \subset D$ и любого $\varepsilon>0$ найдется многочлен $p(z)$ с комплексными коәффициентами такой, что $|f(z)-p(z)|<\varepsilon$ для всех $z \in K$.



В иной формулировке: любая функция $f$, голоморфная в односвязной области $D \subset \mathbb{C}$, представляется в виде ряда из многочленов от $z$, абсолютно и равномерно сходящегося к $f$ внутри $D$. Эквивалентная формулировка Р. т.: пусть $K-$ компакт на комплексной плоскости $\mathbb{C}$ со связным дополнением $\mathbb{C} \backslash K$; тогда всякая функция, голоморфная в окрестности $K$, равномерно на $K$ приближается многочленами от $z$. В такой форме Р. т. есть частный случай Мергеляна теоремы.



Р. т. наз. также следующая т е о р е м а о р а ц и она льны х п р и б л и же ни я х: всякая функция $f$, голоморфная в области $D \subset \mathbb{C}$, равномерно внутри $D$ приближается рациональными функциями с полюсами вне $D$.



Р. т. имеет многочисленные применения в теории функций комплексного переменного и в функциональном анализе. Аналог Р. т. справедлив на некомпактных римановых поверхностях. Обобщением Р. т. для функций многих комплексных переменных является теорема Ока-Вейля (см. Ока теоремы).



Лит.: [1] М а р к у ш е в и ч А. И., Краткий курс теории аналитических функций, 4 изд., М., 1978; [2] III а б а т Б. В. Введение в комплексный анализ, 2 изд., ч. 1, М., 1976.



РУНГЕ - КУТТА МЕТОД - одношаговый метод численного решения задачи Коти для системы обыкновенных диффференциальных уравнений вида



\[

u^{\prime}=f(t, u)

\]



Основная идея Р.- К. м. была предложена К. Рунге [1] и развита затем В. Кутта [2] и др. Первоначально әта идея использовалась лишь для построения явных схем Р.- К. м., к-рые разыскивались в виде



\[

y_{j+1}=y_{j}+\tau \sum_{i=1}^{q} A_{i} k_{i}

\]



где



\[

y_{j+\alpha} \cong u\left(t_{j}+\alpha \tau\right), \quad k_{1}=f\left(t_{j}, y_{j}\right),

\]



$k_{n}=f\left(t_{j}+\alpha_{n} \tau, y_{j}+\tau \sum_{m=1}^{n-1} \beta_{n m} k_{m}\right), n=2,3, \ldots, q$,



при этом значения постоянных $A_{i}, \alpha_{n}, \beta_{n m}, i=1,2, \ldots, q ;$ $n=2,3, \ldots, q ; m=1,2, \ldots, n-1$, определялись из требования, чтобы погрешность равенства (2) на точном решении уравнения (1) имела возможно высокий порядок малости в сравнении с шагом $\tau$ для любых уравнений вида (1).



\includegraphics[max width=\textwidth, center]{2023_07_08_9750cb79bde05e1dc0dbg-1}

тодов, Р.- К. м., как и всякий одношаговый метод, не требует предварительного построения начала таблицы значений приближенного решения и дает возможность вести вычислительный процесс при естественных для уравнения (1) начальных условиях, что позволяет использовать его непосредственно и в случае неравномерных сеток. Однако поскольку в этом методе не используется информация о решении в предыдущих узлах сетки, то он, вообще говоря, оказывается локально менее экономичным, чем, напр., метод Адамса.



Наиболее широко известным (см., напр., [3]) среди Р. - К. м. является метод



\[

\begin{gathered}

y_{j+1}=y_{j}+\frac{1}{6} \tau\left(k_{1}+2 k_{2}+2 k_{3}+k_{4}\right), \\

k_{1}=f\left(t_{j}, y_{j}\right), k_{2}=f\left(t_{j}+\frac{1}{2} \tau, y_{j}+\frac{1}{2} \tau k_{1}\right), \\

k_{3}=f\left(t_{j}+\frac{1}{2} \tau, y_{j}+\frac{1}{2} \tau k_{2}\right), \quad k_{4}=f\left(t_{j}+\tau, y_{j}+\tau k_{3}\right),

\end{gathered}

\]



принадлежащий зависящему от двух свободных параметров семейству методов четвертого порядка точности вида (2) с $q=4$. Популярен и простейший явный Р.К. м. первого порядка точности, получающийся из (2) при $q=1$. Этот метод известен под названием м е т од а Э й л е р а. При значениях $q$, равных 2 и 3 , из (2) могут быть найдены семейства Р.- К. м. второго и третьего порядка точности, зависящие от одного и двух





\end{document}